\section{Related Work}
\label{sec:related}

Shortest path routing literature covers static precomputation and dynamic query-time acceleration. 
Static methods, including Contraction Hierarchies, achieve sub-millisecond query times by adding shortcut edges to the graph. 
These methods require a static network topology. 
Dynamic environments with real-time congestion updates render static shortcut validation obsolete. 

Algorithmic query-time acceleration methods restrict the explored search space. 
The A* search uses Euclidean distance heuristics to guide the search toward the destination. 
Goldberg and Harrelson \cite{goldberg2004alt} present the ALT algorithm, which leverages landmark nodes and the triangle inequality to construct tighter lower bounds. 
Landmark selection algorithms, studied by Fuchs \cite{fuchs2010preprocessing}, determine lower bound quality and search space reduction. 

Priority queue selection dictates the step latency constant $t_{\text{iter}}$. 
Cherkassky et al. \cite{cherkassky1994shortest} evaluate priority queues on sparse networks. 
Their experiments show that Fibonacci heaps suffer from high constant factors and poor cache locality. 
Simple binary heaps outperform Fibonacci heaps on real-world graphs. 
Ahuja et al. \cite{ahuja1990faster} introduce radix heaps to exploit integer edge weights, bounds, and monotonicity. 
Radix heaps achieve near-constant time complexities by grouping keys into dynamic buckets. 

Arthur et al. \cite{quickheap_arxiv} propose QuickHeap to minimize element comparison counts through lazy sorting and deferred element positioning. 
While QuickHeap exhibits strong theoretical bounds, the lazy deferral mechanism introduces branch instruction overhead. 
We examine how these instruction-level overheads degrade CPU performance on metropolitan road networks. 
We contrast these priority queues against our cache-conscious SIMD $d$-ary heap layout.
